% =====================================================================
%  Máxima verossimilhança com dados de RAD-seq no IQ-TREE 3
%  Métodos Computacionais em Sistemática e Filogeografia (UNESP)
%
%  Compilar com pdfLaTeX (Overleaf: Menu > Compiler > pdfLaTeX).
%
%  VERSÃO: por padrão gera a versão do ALUNO (sem as notas do professor).
%  Para a versão do PROFESSOR, troque \professorfalse por \professortrue
%  na linha indicada abaixo, ou compile com
%     pdflatex "\def\versaoprofessor{}\input{apostila_iqtree.tex}"
% =====================================================================
\documentclass[11pt,a4paper]{article}

\newif\ifprofessor
\professorfalse            % <-- troque para \professortrue para incluir as notas do professor
\ifdefined\versaoprofessor \professortrue \fi

\usepackage[T1]{fontenc}
\usepackage[utf8]{inputenc}
\usepackage{lmodern}
\usepackage{textcomp}
\usepackage{amsmath}
\usepackage[provide=*,brazilian]{babel}
\usepackage[a4paper,margin=2.5cm]{geometry}
\usepackage{fancyvrb}
\usepackage{upquote}
\usepackage{booktabs}
\usepackage{array}
\usepackage{enumitem}
\usepackage{comment}
\usepackage{needspace}
\usepackage{url}
\usepackage[hidelinks]{hyperref}

\fvset{xleftmargin=1.5em,fontsize=\small}
\setlist{itemsep=0.2em,topsep=0.4em}
\setlength{\parskip}{0.35em}
\setlength{\emergencystretch}{3em}
\urlstyle{tt}

\newcommand{\arq}[1]{\path{#1}}
\newcommand{\ui}[1]{\textit{#1}}
\newcommand{\cmd}[1]{\texttt{#1}}

\ifprofessor
  \specialcomment{resposta}{\begin{quote}\small\textbf{Para o professor.}\ }{\end{quote}}
\else
  \excludecomment{resposta}
\fi

\newcounter{questao}
\newcommand{\questao}{\par\needspace{3\baselineskip}\refstepcounter{questao}\noindent\textbf{Questão~\thequestao.}\ }

\title{Máxima verossimilhança com dados de RAD-seq no IQ-TREE 3\\[0.3em]
\large Matriz completa, SNPs e fatores de concordância em \emph{Pedicularis}}
\author{Métodos Computacionais em Sistemática e Filogeografia\\Universidade Estadual Paulista (UNESP)}
\date{\ifprofessor Versão do professor (com notas)\else Versão do aluno\fi\\ Novembro de 2026}

\begin{document}
\maketitle

\paragraph{O roteiro.} Vamos estimar a filogenia das 13 amostras de \emph{Pedicularis} por
máxima verossimilhança no IQ-TREE 3 \cite{wong2025}, com três matrizes que o ipyrad gerou a partir
dos mesmos dados, e comparar o que cada uma mostra. No fim, vamos medir o conflito entre os sítios
com os fatores de concordância de sítio.

\paragraph{Arquivos.} A pasta \arq{pratica_iqtree} contém as três matrizes do ipyrad, para quem
não tiver as suas: \arq{pedicularis.phy}, \arq{pedicularis.snps} e \arq{pedicularis.usnps}.

\paragraph{Programas.} IQ-TREE 3 e FigTree.

\begin{resposta}
Os números do gabarito foram obtidos com o IQ-TREE 3.1.4 (Linux) e as matrizes do ipyrad 0.9.108
que você gerou. Quem usar as próprias saídas do ipyrad vai ter números um pouco diferentes; com
\cmd{-{}-seed 2026} e os arquivos da pasta, devem ser os mesmos do gabarito. Na máquina dos testes, com
dois núcleos, os tempos foram: Análise~1, 17 minutos; Análise~2, 10 minutos (com um núcleo);
Análise~3, 2,5 minutos; e o sCF, 15 minutos (quase todo gasto reestimando os parâmetros do modelo
na matriz completa). Num notebook com mais núcleos, deve ser mais rápido. Vale começar pela
Análise~1 e deixá-la rodando enquanto os alunos fazem as Análises 2 e 3 em outra janela do
terminal; o sCF entra logo depois que ela terminar.
\end{resposta}

% =====================================================================
\section{Antes de começar: o tutorial do ipyrad}
% =====================================================================

Esta prática parte das saídas do tutorial do ipyrad com os dados de RAD-seq de \emph{Pedicularis}
(13 amostras, Eaton \& Ree \cite{eaton2013}), feito na aula anterior:
\url{https://ipyrad.readthedocs.io/en/master/pedicularis_.html}. Três avisos sobre aquele
tutorial:
\begin{enumerate}
\item \textbf{O download pode falhar.} O \cmd{ipyrad -{}-download} usa o \cmd{fastq-dump} do
sra-tools, e o ambiente do ipyrad pode vir com uma versão de 2019 (2.9.6), que não lê o formato
atual dos arquivos do NCBI. O erro aparece como \cmd{rcColumn, rcNotFound}. Para consertar,
dentro do ambiente do ipyrad:
\begin{Verbatim}
conda install -c conda-forge -c bioconda "sra-tools>=3"
fastq-dump --version
\end{Verbatim}
A versão deve começar com 3. Apague os arquivos \arq{.sra} que tiverem sobrado na pasta e rode
o \cmd{ipyrad -{}-download} de novo.
\item \textbf{O arquivo de SNPs mudou de nome.} Quando o tutorial manda olhar o
\arq{pedicularis.snps.phy}, na versão atual do ipyrad o arquivo é o \arq{pedicularis.snps} (o
mesmo formato PHYLIP, só sem o \arq{.phy} no nome).
\item \textbf{Cuidado com os parâmetros 25 e 26.} O tutorial muda o 26 (\cmd{trim\_loci}) para
\cmd{0, 5, 0, 0}. O 25 (\cmd{trim\_reads}), logo acima, tem o mesmo formato: se o valor for parar
nele, cada sequência é cortada para 5 pb e o passo 2 descarta todas
(\cmd{No reads passed filtering}).
\end{enumerate}

% =====================================================================
\section{Introdução}
% =====================================================================

\paragraph{As três matrizes.} O ipyrad escreve a mesma informação de vários jeitos. Vamos usar
três:
\begin{center}\small
\begin{tabular}{@{}lp{7.6cm}c@{}}
\toprule
Arquivo & O que tem & Correção +ASC?\\
\midrule
\arq{pedicularis.phy} & todos os sítios de todos os loci, concatenados & não\\
\arq{pedicularis.snps} & só os sítios variáveis: todos os SNPs de cada locus & sim\\
\arq{pedicularis.usnps} & um SNP sorteado por locus (\emph{unlinked}) & sim\\
\bottomrule
\end{tabular}
\end{center}

\paragraph{Por que corrigir.} Uma matriz só de SNPs não tem sítios invariáveis, porque eles foram
tirados. Se o modelo não souber disso, ele interpreta ``todo sítio variou'' como uma taxa de
substituição altíssima, e os ramos saem longos demais. A correção do viés de averiguação
(\emph{ascertainment bias correction}, \cmd{+ASC} no IQ-TREE \cite{lewis2001}) calcula a
verossimilhança de cada sítio \emph{sabendo que ele é variável}: divide a verossimilhança do sítio
pela probabilidade de um sítio qualquer ser variável,
\[
L_i^{\text{ASC}} = \frac{L_i}{1 - \left[L(\text{AAA}\ldots) + L(\text{CCC}\ldots) +
L(\text{GGG}\ldots) + L(\text{TTT}\ldots)\right]},
\]
em que $L(\text{AAA}\ldots)$ é a verossimilhança de uma coluna com A em todas as amostras.

\paragraph{Por que olhar além do \emph{bootstrap}.} Com centenas de milhares de sítios, o
\emph{bootstrap} ultrarrápido (UFBoot) \cite{hoang2018} tende a dar 100 em quase todo nó, mesmo
onde há conflito real entre os sítios (por separação incompleta de linhagens, por exemplo). O
fator de concordância de sítio (sCF) \cite{minh2020,mo2023} mede outra coisa: a porcentagem dos
sítios informativos que apoiam cada ramo, contra as duas alternativas. Vamos calcular os dois.

% =====================================================================
\section{Instalação e arquivos}
% =====================================================================

\paragraph{O IQ-TREE 3.} Os executáveis estão em \url{https://iqtree.github.io}, e o programa se
chama \cmd{iqtree3}. Use a versão 3.1.4 ou mais nova: foi a dos testes desta apostila. Como as saídas do ipyrad estão no Linux (ou no WSL) ou no Mac, instale-o
no mesmo lugar: baixe o arquivo do seu sistema (por exemplo,
\arq{iqtree-3.1.4-Linux-intel.tar.gz} no Linux e no WSL) e descompacte com
\cmd{tar xzf} seguido do nome do arquivo. O programa é o \arq{iqtree3} da pasta \arq{bin}.

\paragraph{O \texttt{PATH}.} Até aqui, rodamos programas baixados com \cmd{./}, de dentro da
pasta onde eles estão. O jeito mais prático é colocá-los numa pasta onde o terminal sempre
procura comandos. Essa lista de pastas é a variável de ambiente \cmd{PATH}; veja quais são com
\begin{Verbatim}
echo $PATH
\end{Verbatim}
A pasta \arq{/usr/local/bin} está nessa lista no Linux, no WSL e no Mac. Mova o \arq{iqtree3}
para lá (o \cmd{sudo} pede a senha de administrador; no WSL, é a senha que você criou ao instalar
o Ubuntu):
\begin{Verbatim}
sudo mkdir -p /usr/local/bin
sudo mv iqtree-3.1.4-Linux-intel/bin/iqtree3 /usr/local/bin/
iqtree3 --version
\end{Verbatim}
O primeiro comando só cria a pasta, se ela ainda não existir (nos Macs com chip Apple, ela às
vezes não existe). No segundo, troque \arq{iqtree-3.1.4-Linux-intel} pelo nome da pasta que
saiu do arquivo do seu sistema. No Mac, antes do \cmd{mv}, rode \cmd{xattr -d com.apple.quarantine} seguido
do caminho do \arq{iqtree3}, para tirar o bloqueio que o Mac põe em programas baixados. Agora o
\cmd{iqtree3} funciona de qualquer pasta.

\paragraph{Os arquivos.} Crie uma pasta para a prática e copie para ela as três matrizes, que
ficam na pasta de saídas do ipyrad (ou use as da pasta da prática):
\begin{Verbatim}
mkdir iqtree
cp analysis-ipyrad/pedicularis_outfiles/pedicularis.phy iqtree/
cp analysis-ipyrad/pedicularis_outfiles/pedicularis.snps iqtree/
cp analysis-ipyrad/pedicularis_outfiles/pedicularis.usnps iqtree/
cd iqtree
\end{Verbatim}

\paragraph{O grupo externo.} As duas amostras de \emph{P. przewalskii} enraízam as árvores. Os
nomes vêm com o número do SRA, e o IQ-TREE recebe os dois separados por vírgula, sem espaço:
\begin{Verbatim}
-o 32082_przewalskii_SRR1754729,33588_przewalskii_SRR1754727
\end{Verbatim}
Para conferir os nomes, olhe o começo de cada linha do \arq{pedicularis.snps} com
\cmd{cut -c1-40 pedicularis.snps}.

\paragraph{Os comandos.} Os comandos abaixo terminam algumas linhas com \cmd{\textbackslash}: é o
jeito de quebrar um comando longo em várias linhas no terminal. Pode copiar assim mesmo, ou
digitar tudo numa linha só, sem as barras.

% =====================================================================
\section{Análise 1: a matriz completa}
% =====================================================================

A primeira árvore usa todos os sítios e serve de referência para as outras duas. O IQ-TREE busca
a árvore de máxima verossimilhança e calcula o UFBoot com 1.000 réplicas (\cmd{-B 1000}):
\begin{Verbatim}
iqtree3 -s pedicularis.phy -st DNA -m GTR+I+G4 -B 1000 -T AUTO \
  -o 32082_przewalskii_SRR1754729,33588_przewalskii_SRR1754727 \
  --prefix completo --seed 2026
\end{Verbatim}
\begin{itemize}
\item \cmd{-st DNA}: avisa que os dados são DNA. Sem isso, o IQ-TREE para com o erro
\cmd{cannot confidently detect the sequence data type}: a matriz tem tantos \cmd{N} (dados
faltantes) e códigos de ambiguidade (os heterozigotos) que ele não tem certeza de que é DNA.
\item \cmd{-m GTR+I+G4}: o modelo de substituição, fixado. Nas Análises 2 e 3, o IQ-TREE vai
escolher o modelo sozinho, com o ModelFinder \cite{kalyaanamoorthy2017}; aqui, com 2,8 milhões
de colunas, a escolha levaria horas. Usamos o modelo mais geral, o GTR, com variação de taxas
entre sítios (\cmd{+G4}) e uma proporção de sítios invariáveis (\cmd{+I}), já que quase todos os
sítios da matriz são constantes.
\item \cmd{-T AUTO}: o IQ-TREE escolhe quantos núcleos do processador usar.
\item \cmd{-{}-prefix}: o começo do nome de todos os arquivos de saída.
\item \cmd{-{}-seed}: a semente dos números aleatórios, para todos chegarem ao mesmo resultado.
\end{itemize}
Esta é a análise mais demorada da prática: deixe-a rodando e, enquanto isso, faça as Análises 2
e 3 em outra janela do terminal. Os principais arquivos de saída:
\begin{itemize}
\item \arq{completo.iqtree}: o relatório, com o modelo escolhido, o logaritmo da verossimilhança,
o comprimento total da árvore e a árvore desenhada em texto;
\item \arq{completo.treefile}: a árvore de máxima verossimilhança, com o UFBoot em cada nó;
\item \arq{completo.log}: tudo o que apareceu na tela.
\end{itemize}
Abra o \arq{completo.treefile} no FigTree, enraíze nos dois \emph{P. przewalskii} e mostre os
rótulos dos nós (\ui{Node Labels}).

\questao No \arq{completo.iqtree}, quantos sítios a matriz tem? Quantos são constantes, quantos
são informativos para parcimônia (\emph{parsimony-informative}) e quantos são
\emph{singletons}? Qual proporção de sítios invariáveis o modelo estimou (\cmd{Proportion of
invariable sites})? Compare com a proporção de sítios constantes.

\begin{resposta}
A matriz tem 2.826.933 sítios: 2.658.382 constantes (94,0\%), 73.991 informativos para
parcimônia e 94.560 \emph{singletons}. A proporção de invariáveis estimada é 0,627, bem menor
que os 94\% de colunas constantes: muitas colunas constantes são só sítios lentos que não
mudaram nestas 13 amostras (o gama, com $\alpha = 0{,}79$, dá conta deles), e boa parte delas
tem dados faltantes, o que as torna constantes com pouca informação.
\end{resposta}

\questao As amostras da mesma espécie formam clados? Quantos nós têm UFBoot abaixo de 95 (o
limite que a documentação do IQ-TREE sugere)?

\begin{resposta}
Com o grupo externo, a árvore tem dois grandes clados: (\emph{superba},
\emph{cyathophylla}) + \emph{cyathophylloides}, e \emph{thamnophila} + \emph{rex}. As duas
\emph{cyathophylloides} formam um clado (100), assim como as cinco \emph{rex} (99).
\textbf{As duas \emph{thamnophila} não formam um clado}: \cmd{30556\_thamno} é irmã de
\emph{rex} (100), e \cmd{33413\_thamno} é irmã das duas (100). \emph{superba} e
\emph{cyathophylla} têm uma amostra cada. Nenhum nó tem UFBoot abaixo de 95; o menor é 96,
dentro de \emph{rex}.
\end{resposta}

% =====================================================================
\section{Análise 2: só os SNPs, com +ASC}
% =====================================================================

Agora a matriz só de SNPs. O \cmd{+ASC} no fim do \cmd{-m MFP} faz o ModelFinder testar só
modelos com a correção:
\begin{Verbatim}
iqtree3 -s pedicularis.snps -st DNA -m MFP+ASC -B 1000 -T AUTO \
  -o 32082_przewalskii_SRR1754729,33588_przewalskii_SRR1754727 \
  --prefix snps --seed 2026
\end{Verbatim}

\paragraph{O erro esperado.} O IQ-TREE vai parar logo no começo, com esta mensagem:
\begin{Verbatim}[fontsize=\footnotesize]
ERROR: For your convenience alignment with variable sites printed to snps.varsites.phy
ERROR: Invalid use of +ASC because of 33197 invariant sites in the alignment
\end{Verbatim}
Mas o ipyrad não tinha tirado os sítios invariáveis? Tinha, só que ``variável'' não quer dizer a
mesma coisa para os dois programas. Para o ipyrad, um sítio onde uma amostra é heterozigota
(\cmd{R}, por exemplo, que é A ou G) já tem um SNP. Para o +ASC, a coluna só é variável se não
puder ser explicada por uma única base: uma coluna só com \cmd{A}, \cmd{R} e \cmd{N} pode ser
toda de A, e conta como invariável.

O IQ-TREE já resolve: junto com o erro, ele grava um alinhamento só com as colunas realmente
variáveis, no arquivo \arq{snps.varsites.phy} (o nome vem do \cmd{-{}-prefix}). Rode de novo
com ele:
\begin{Verbatim}
iqtree3 -s snps.varsites.phy -st DNA -m MFP+ASC -B 1000 -T AUTO \
  -o 32082_przewalskii_SRR1754729,33588_przewalskii_SRR1754727 \
  --prefix snps --seed 2026 -redo
\end{Verbatim}
O \cmd{-redo} é importante. A tentativa que deu erro deixou um arquivo de checkpoint
(\arq{snps.model.gz}), e, sem o \cmd{-redo}, o IQ-TREE tenta aproveitá-lo, mesmo ele sendo de
outro alinhamento. Dependendo da versão, isso faz o programa travar no meio do ModelFinder, com
um erro \cmd{Assertion ... failed} e \cmd{IQ-TREE CRASHES WITH SIGNAL ABORTED}. O \cmd{-redo}
manda ignorar o checkpoint e começar do zero.

\questao Quantas colunas o IQ-TREE tirou? Que fração das colunas do \arq{pedicularis.snps} isso
representa? Compare o número de colunas que sobraram com o número de sítios variáveis da
Análise 1 (informativos mais \emph{singletons}).

\begin{resposta}
O IQ-TREE tirou 33.197 colunas das 201.748 (16,5\%). Sobraram 168.551, que são
exatamente os 73.991 informativos mais os 94.560 \emph{singletons} da Análise~1: os mesmos
sítios variáveis, vistos pelos dois lados.
\end{resposta}

\questao Qual modelo o ModelFinder escolheu, e por qual critério? Compare o comprimento total da
árvore (no \arq{.iqtree}) com o da Análise 1. Por que eles são tão diferentes, se os dados são
os mesmos?

\begin{resposta}
O ModelFinder escolheu o TVM+ASC+G4, pelo BIC (é o critério padrão do IQ-TREE), quase
empatado com o TVMe+ASC+G4 (pesos de BIC de 0,33 e 0,32). O comprimento total da árvore é de
0,560 substituição por sítio, contra 0,117 na Análise~1, quase cinco vezes mais. Na Análise~1,
o comprimento é médio sobre 2,8 milhões de sítios, 94\% deles constantes. Com o +ASC, o modelo
só sabe que os sítios constantes foram tirados, não quantos: ele supõe que eram só os que o
próprio modelo (com o gama) espera, que são muito menos do que os que existiam de verdade. Por
isso, os comprimentos de ramo de uma matriz de SNPs não se comparam diretamente com os de uma
matriz completa; as proporções entre os ramos, sim.
\end{resposta}

\questao (Opcional.) Rode a Análise 2 \emph{sem} a correção, com o modelo que o ModelFinder
escolheu, mas sem o \cmd{+ASC} (para não esperar o ModelFinder de novo):
\begin{Verbatim}
iqtree3 -s snps.varsites.phy -st DNA -m TVM+G4 -B 1000 -T AUTO \
  -o 32082_przewalskii_SRR1754729,33588_przewalskii_SRR1754727 \
  --prefix snps_sem_asc --seed 2026
\end{Verbatim}
(Se o seu ModelFinder tiver escolhido outro modelo, use o seu, sem o \cmd{+ASC}.) O que acontece
com o comprimento total da árvore? E com a topologia?

\begin{resposta}
Sem a correção, o comprimento total da árvore vai para 2,30 substituições por sítio, contra
0,560 com o +ASC (quatro vezes mais) e 0,117 na matriz completa (20 vezes mais): sem saber que os
sítios constantes foram tirados, o modelo explica ``todo sítio variou'' com taxas altíssimas. A
topologia é a mesma (RF = 0), mas o UFBoot cai em três nós (\emph{rex}, 87;
\cmd{30556\_thamno} + \emph{rex}, 76; \cmd{35236\_rex} + (\cmd{38362\_rex}, \cmd{39618\_rex}),
83): um modelo errado também afeta o suporte.
\end{resposta}

% =====================================================================
\section{Análise 3: um SNP por locus}
% =====================================================================

Os SNPs de um mesmo locus estão a poucas dezenas de pares de bases uns dos outros e, quase
sempre, compartilham a mesma história genealógica. Na matriz de todos os SNPs, um locus com dez
SNPs pesa dez vezes mais que um locus com um. O \arq{pedicularis.usnps} guarda um SNP sorteado
por locus, mais perto da ideia de sítios independentes que o modelo assume:
\begin{Verbatim}
iqtree3 -s pedicularis.usnps -st DNA -m MFP+ASC -B 1000 -T AUTO \
  -o 32082_przewalskii_SRR1754729,33588_przewalskii_SRR1754727 \
  --prefix usnps --seed 2026
\end{Verbatim}
O mesmo erro do +ASC vai aparecer: rode de novo com o \arq{usnps.varsites.phy}, como na Análise 2,
sem esquecer o \cmd{-redo}:
\begin{Verbatim}
iqtree3 -s usnps.varsites.phy -st DNA -m MFP+ASC -B 1000 -T AUTO \
  -o 32082_przewalskii_SRR1754729,33588_przewalskii_SRR1754727 \
  --prefix usnps --seed 2026 -redo
\end{Verbatim}

\questao Quantas colunas tem a matriz final, comparada à da Análise 2? O que aconteceu com os
valores de UFBoot? A topologia mudou?

\begin{resposta}
A matriz final tem 31.992 colunas (o IQ-TREE tirou 6.960 das 38.952), cerca de um
quinto das 168.551 da Análise~2. O UFBoot caiu em dois nós de \emph{rex}: o clado \emph{rex}
foi de 100 para 97, e um nó novo, que junta \cmd{35236\_rex} a (\cmd{35855\_rex},
\cmd{40578\_rex}), tem só 54. A topologia mudou nesse ponto: nas Análises 1 e 2,
\cmd{35236\_rex} fica com (\cmd{38362\_rex}, \cmd{39618\_rex}), com 96.
\end{resposta}

% =====================================================================
\section{Comparando as árvores}
% =====================================================================

\paragraph{As topologias.} A distância de Robinson-Foulds (RF) conta quantos clados estão numa
árvore e não na outra; 0 quer dizer a mesma topologia. Junte as três árvores num arquivo e peça
ao IQ-TREE as distâncias entre todos os pares:
\begin{Verbatim}
cat completo.treefile snps.treefile usnps.treefile > tres_arvores.tre
iqtree3 -t tres_arvores.tre -rf_all
\end{Verbatim}
As distâncias ficam no \arq{tres_arvores.tre.rfdist}, numa matriz em que as árvores aparecem
na ordem do arquivo (\cmd{Tree0} é a da matriz completa, \cmd{Tree1} a dos SNPs e \cmd{Tree2} a
de um SNP por locus).

\paragraph{Os números.} Preencha com os relatórios (\arq{.iqtree}) das três análises:
\begin{center}\small
\begin{tabular}{@{}lccc@{}}
\toprule
& Completa & Todos os SNPs & Um SNP por locus\\
\midrule
Colunas analisadas & & &\\
Modelo & GTR+I+G4 & &\\
Comprimento total da árvore & & &\\
Nós com UFBoot $<$ 95 & & &\\
Distância RF até a completa & 0 & &\\
\bottomrule
\end{tabular}
\end{center}

\questao Onde estão as diferenças entre as três árvores? Elas aparecem em nós com UFBoot alto ou
baixo?

\begin{resposta}
As Análises 1 e 2 dão a mesma topologia (RF = 0); a Análise~3 difere das duas em um
clado (RF = 2: um clado presente numa árvore e ausente na outra conta duas vezes, uma para cada
lado). A única diferença é a posição de \cmd{35236\_rex}, dentro de \emph{rex}, justamente no
nó com o menor UFBoot das Análises 1 e 2 (96) e no de 54 da Análise~3.
\end{resposta}

% =====================================================================
\section{O fator de concordância de sítio}
% =====================================================================

Na árvore da matriz completa, vamos calcular o sCF de cada ramo. Para cada ramo, o IQ-TREE
sorteia 100 quartetos de amostras em volta dele e, em cada quarteto, vê quais sítios apoiam o
ramo e quais apoiam cada uma das duas alternativas. A árvore (\cmd{-te}) e o modelo são os da
Análise 1:
\begin{Verbatim}
iqtree3 -te completo.treefile -s pedicularis.phy -st DNA -m GTR+I+G4 \
  --scfl 100 -T AUTO --prefix concord --seed 2026
\end{Verbatim}
\begin{itemize}
\item \arq{concord.cf.tree}: a árvore com o UFBoot e o sCF em cada nó, separados por uma barra
(por exemplo, \cmd{100/44.6});
\item \arq{concord.cf.stat}: uma tabela com uma linha por ramo, com o sCF e os dois fatores de
discordância, \cmd{sDF1} e \cmd{sDF2} (a porcentagem dos sítios que apoiam cada alternativa), e o
número médio de sítios informativos por quarteto (\cmd{sN});
\item \arq{concord.cf.branch}: a mesma árvore, com o número de cada ramo (a coluna \cmd{ID} do
\arq{concord.cf.stat}) no lugar dos suportes.
\end{itemize}
Abra o \arq{concord.cf.tree} e o \arq{concord.cf.branch} no FigTree, lado a lado, para saber a que
ramo corresponde cada linha da tabela. (A árvore sai sem raiz: enraíze nos dois \emph{P.
przewalskii}, como antes.)

\paragraph{Como ler.} Em volta de cada ramo, há três arranjos possíveis dos quatro grupos de
amostras. Um sCF perto de 33\% quer dizer que os sítios não preferem o ramo da árvore a nenhuma
das alternativas; um sCF bem acima disso quer dizer que a maioria dos sítios concorda com ele.
Um ramo pode ter UFBoot de 100 e sCF de 40\%: é o mais apoiado dos três arranjos, mas com muito
conflito. Os dois fatores de discordância também contam uma história: se sDF1 e sDF2 são
parecidos, o conflito é o esperado pela separação incompleta de linhagens; se um é bem maior que
o outro, um dos arranjos alternativos tem apoio demais, e isso pode indicar fluxo gênico entre
as linhagens desse arranjo.

\questao Quais ramos têm UFBoot alto (95 ou mais) e sCF abaixo de 50\%? Eles estão entre
espécies ou entre amostras da mesma espécie?

\begin{resposta}
Três ramos: \emph{thamnophila} 30556 + \emph{rex} (UFBoot 100, sCF 44,6\%), o clado
\emph{rex} (99; 37,9\%) e \cmd{35236\_rex} + (\cmd{38362\_rex}, \cmd{39618\_rex}) (96;
35,0\%). Os dois primeiros estão entre espécies (\emph{thamnophila} e \emph{rex}); o terceiro,
dentro de \emph{rex}. Outros ramos com UFBoot de 100 também têm sCF modesto: (\emph{superba},
\emph{cyathophylla}), 65,6\%; e esse par mais \emph{cyathophylloides}, 65,4\%. Só os ramos
dentro de \emph{cyathophylloides} (99,9\%), o do grupo externo (99,7\%) e o de (\cmd{38362\_rex},
\cmd{39618\_rex}) (96,1\%) têm sCF perto de 100.
\end{resposta}

\questao Para esses ramos, sDF1 e sDF2 são parecidos ou um é bem maior que o outro? O que isso
sugere?

\begin{resposta}
No clado \emph{rex} + \cmd{30556\_thamno} (ramo 17 no \arq{concord.cf.branch}),
sDF1 = 36,0\% e sDF2 = 19,4\%: uma das alternativas tem quase o dobro do apoio da outra, o que
sugere fluxo gênico entre as linhagens desse arranjo (ou algum outro viés). No clado \emph{rex}
(ramo 18), 26,4\% e 35,7\%: também assimétrico, e a alternativa mais apoiada fica perto do
próprio sCF (37,9\%). No ramo 19, dentro de \emph{rex}, 33,4\% e 31,6\%, com sCF de 35,0\%: os
três arranjos têm praticamente o mesmo apoio, o esperado para um ramo muito curto, com separação
incompleta de linhagens e pouco sinal. Bom lembrar que \emph{thamnophila} e \emph{rex} são
representadas por poucas amostras: o padrão é uma pista, não uma conclusão.
\end{resposta}

% =====================================================================
\section{Na vida real}
% =====================================================================

\paragraph{Um segundo suporte.} O UFBoot é rápido e, com modelos adequados, bem calibrado: um
clado com 95 é verdadeiro em cerca de 95\% dos casos. Mas ele pode superestimar o apoio quando o
modelo está errado. Por isso, a documentação do IQ-TREE recomenda calcular junto o SH-aLRT
(\cmd{-alrt 1000}) \cite{guindon2010}, um teste que compara cada ramo com os dois arranjos
alternativos em volta dele, reamostrando os sítios. Os dois saem no mesmo rótulo do nó
(SH-aLRT/UFBoot), e um clado é considerado confiável quando o SH-aLRT é pelo menos 80 e o UFBoot
pelo menos 95. Numa análise para publicação, rodaríamos sempre os dois.

\paragraph{A escolha do modelo.} Na Análise~1, fixamos o GTR+I+G4 porque o ModelFinder levaria
horas numa matriz de 2,8 milhões de colunas. Numa análise para publicação, escolheríamos o modelo
mesmo assim, deixando o computador trabalhar ou limitando os modelos testados (por exemplo, só
variação de taxas por gama, com \cmd{-mrate G,I+G}).

\paragraph{Além da concatenação.} Uma árvore de máxima verossimilhança de dados concatenados
assume que todos os loci têm a mesma história. Com RAD-seq de espécies próximas, isso raramente
é verdade, e o sCF mostra o tamanho do conflito. Numa análise para publicação, complementaríamos
com métodos de árvore de espécies que lidam com essa discordância, como o SVDquartets ou o
tetrad (das ferramentas de análise do ipyrad). O fator de concordância de gene (gCF) exigiria
uma árvore por locus, e os loci de RAD-seq (cerca de 70 pb aqui) são curtos demais para dar
árvores informativas.

% =====================================================================
\begin{thebibliography}{99}
\bibitem{eaton2013} Eaton, D.~A.~R. \& Ree, R.~H. 2013. Inferring phylogeny and introgression
using RADseq data: an example from flowering plants (\emph{Pedicularis}: Orobanchaceae).
\emph{Systematic Biology} 62: 689--706.
\bibitem{guindon2010} Guindon, S., Dufayard, J.-F., Lefort, V., Anisimova, M., Hordijk, W. \&
Gascuel, O. 2010. New algorithms and methods to estimate maximum-likelihood phylogenies:
assessing the performance of PhyML 3.0. \emph{Systematic Biology} 59: 307--321.
\bibitem{hoang2018} Hoang, D.~T., Chernomor, O., von Haeseler, A., Minh, B.~Q. \& Vinh, L.~S.
2018. UFBoot2: improving the ultrafast bootstrap approximation. \emph{Molecular Biology and
Evolution} 35: 518--522.
\bibitem{kalyaanamoorthy2017} Kalyaanamoorthy, S., Minh, B.~Q., Wong, T.~K.~F., von Haeseler, A.
\& Jermiin, L.~S. 2017. ModelFinder: fast model selection for accurate phylogenetic estimates.
\emph{Nature Methods} 14: 587--589.
\bibitem{lewis2001} Lewis, P.~O. 2001. A likelihood approach to estimating phylogeny from
discrete morphological character data. \emph{Systematic Biology} 50: 913--925.
\bibitem{minh2020} Minh, B.~Q., Hahn, M.~W. \& Lanfear, R. 2020. New methods to calculate
concordance factors for phylogenomic datasets. \emph{Molecular Biology and Evolution} 37:
2727--2733.
\bibitem{mo2023} Mo, Y.~K., Lanfear, R., Hahn, M.~W. \& Minh, B.~Q. 2023. Updated site
concordance factors minimize effects of homoplasy and taxon sampling. \emph{Bioinformatics} 39:
btac741.
\bibitem{wong2025} Wong, T.~K.~F., Ly-Trong, N., Ren, H. \emph{et al.} 2025. IQ-TREE 3:
phylogenomic inference software using complex evolutionary models. Preprint, EcoEvoRxiv,
doi:10.32942/X2P62N.
\end{thebibliography}

\end{document}
